\subsection{Μηχανισμοί παραγωγής θερμότητας}
\label{subsec:heat_gen}

Το σύνολο της θερμότητας που παράγεται στον τροχό προέρχεται κυρίως από το ελαστικό
και το σύστημα πέδησης. Ενώ η θερμότητα του συστήματος πέδησης μπορεί να επηρεάσει σημαντικά
τον όγκο ελέγχου του τροχού, στην παρούσα εργασία θα εστιάσουμε στην παραγωγή θερμότητας από
τα ελαστικά θεωρώντας πως δεν υπάρχει συναλλαγή με το σύστημα πέδησης, άλλωστε σε αρκετές
έρευνες έχει χρησιμοποιηθεί αυτή η θεώρηση με αξιοσημείωτη ακρίβεια \cite{West2022, Tremlett2016}.

Η ενέργεια που εισέρχεται στο ελαστικό κατά τη λειτουργία προέρχεται από δύο ξεχωριστές
φυσικές διεργασίες. Η πρώτη και κυρίαρχη σε έντονες μανούβρες είναι η \emph{τριβή
ολίσθησης} (sliding friction) στο πέλμα: όταν υπάρχει σχετική κίνηση μεταξύ
ελαστικού και οδοστρώματος, το μηχανικό έργο της τριβής μετατρέπεται αδιαβατικά σε
θερμότητα. Ο ρυθμός αυτής της παραγωγής εξαρτάται τόσο από την ένταση των δυνάμεων στο
πέλμα (εφαπτομενικές δυνάμεις $F_x$, $F_y$) όσο και από τον βαθμό ολίσθησης, που
χαρακτηρίζεται από τον λόγο ολίσθησης $\kappa$ και τη γωνία ολίσθησης $\alpha$.

Η δεύτερη πηγή είναι η \emph{υστέρηση} (hysteresis) του σκελετού του ελαστικού: το ελαστομερές δεν
συμπεριφέρεται ως ιδανικά ελαστικό υλικό αλλά ως ιξωδοελαστικό -- κατά την κυκλική φόρτισή
του (κάθε στροφή του τροχού) ένα μέρος της παραμορφωτικής ενέργειας δεν αποκαθίσταται αλλά
αποδίδεται ως θερμότητα. Αυτός ο μηχανισμός είναι σημαντικός ακόμα και απουσία ολίσθησης.

Αξίζει να σημειωθεί η συμβολή του κατακόρυφου φορτίου $F_z$: η κάμψη του πλευρικού τοιχώματος
-- εντονότερη υπό μεγαλύτερο φορτίο -- συνεισφέρει σημαντικά στην παραγωγή θερμότητας, 
ακόμη και χωρίς ιδιαίτερα υψηλές πλευρικές και διαμήκεις επιταχύνσεις.
Ο Farroni κ.α. \cite{Farroni2014TRT} αναλύουν αυτόν
τον μηχανισμό λεπτομερώς στο πλαίσιο ενός τρισδιάστατου θερμικού
μοντέλου πεπερασμένων στοιχείων.

% ------------------------------------------------------------------
\subsection{Μηχανισμοί απαγωγής θερμότητας}
\label{subsec:heat_diss}
% ------------------------------------------------------------------
Η θερμοκρασία μόνιμης κατάστασης του πέλματος καθορίζεται από τη θερμότητα που παράγεται
αλλά και από εκείνη που απάγεται. Δύο κύριοι μηχανισμοί απαγωγής δρουν ταυτόχρονα κατά τη
λειτουργία.

Ο πρώτος είναι η \emph{συναγωγή με τον αέρα} (convective cooling) λόγω της έκθεσης του ελαστικού
σε ροή αέρα εξωτερικά. Η ψύξη αυτή εξαρτάται από τη διαφορά θερμοκρασίας
επιφάνειας–αέρα και από τη σχετική ταχύτητα. Η γενική εξίσωση συναγωγής φαίνεται στην \cref{eq:convection},
ο συντελεστής συναγωγής είναι τυπικά της μορφής $h = \alpha V^n$ επομένως επηρεάζεται κυρίως
απο την ταχύτητα του οχήματος και είναι ο βασικότερος μηχανισμός που το ελαστικό απάγει θερμότητα.
Ωστόσο, ειδικά σε αγωνιστικές εφαρμογές, δεν είναι πάντα λύση να εκτεθεί στη ροή του αέρα (ή να προστεθέι
αεραγωγός ψύξης προς το ελαστικό). Τα αγωνιστικά οχήματα είναι πολύ ευαίσθητα στον ομόρρου που δημιουργείται
απο τα ελαστικά ο οποίος αυξάνει την αεροδυναμική αντίσταση \cite{tyreWake}, οπότε στην πράξη οι σχεδιαστές προτιμούν
να περιορίσουν την έκθεση του ελαστικού στη ροή, και εφόσον χρειαστεί η θερμοκρασία των ελαστικών ελαττώνεται 
περιορίζοντας την παραγωγή θερμότητας. 


\begin{equation}
    Q_{\text{συναγωγή}} = hA(T_{\text{πέλμα}} - T_{\text{αέρας}})
    \label{eq:convection}
\end{equation}


Ο δεύτερος μηχανισμός είναι η \emph{αγωγιμότητα με το οδόστρωμα} (tyre-road conduction).
Στο πέλμα, όπου το ελαστικό πιέζεται πάνω στην άσφαλτο, υπάρχει άμεση ανταλλαγή
θερμότητας \cref{eq:conduction}.

\begin{equation}
    Q_{αγωγή} = \dfrac{T_{\text{πέλμα}} - T_{\text{οδοστ.}}}{R_{\text{επαφής}}}
    \label{eq:conduction}
\end{equation}

Σημειώνεται εδώ πως η θερμότητα $Q_{\text{αγωγή}}$ μπορεί να είναι αρνητική -- δηλαδή ο
δρόμος να \emph{θερμαίνει} το ελαστικό -- ειδικά κατά τους πρώτους γύρους με κρύα
ελαστικά, εφόσον $T_{\text{οδόστρομα}} > T_{\text{πέλμα}}$. Αυτό
συμβαίνει ιδιαίτερα λόγω του μηχανισμού αγωγής, διότι συνήθως (και ειδικότερα
το καλοκαίρι ή το μεσημέρι) το οδόστρομα είναι αρκετούς βαθμούς θερμότερο
από τη θερμοκρασία περιβάλλοντος.
